%% ma: L11-B04-0023
%% lop: 11
%% chuong: 1
%% bai: 4
%% dang: 11.4.12
%% loai: DS
%% muc: [NB, NB, NB, VD]
%% dap_an: "SSSĐ"
%% nguon: "de-on-tap-giua-ki-1-toan-11-nam-2026-2027-truong-thpt-ngo-quyen-da-nang.pdf (M09, Đề ôn tập giữa kì 1 Toán 11 THPT Ngô Quyền, Đà Nẵng), Mã đề 0113, Phần II Câu 1"
%% kien_thuc: "Bài 2 (công thức lượng giác), Bài 3 (hàm số lượng giác), Bài 4 (phương trình lượng giác cơ bản)"
%% trang_thai: cho-duyet
%% ly_do: "Dạng mới đề xuất 11.4.12"
\begin{ex}
Cho hàm số $f(x)=\sin2x$. Xét tính đúng sai của các mệnh đề sau:
\choiceTF{Hàm số $f(x)$ là hàm số chẵn và tuần hoàn với chu kì $\pi$}{Tập xác định của hàm số $f(x)$ là $D=\mathbb R\setminus\{k\pi,\ k\in\mathbb Z\}$}{Công thức biến đổi $f(x)=\sin x\cos x$}{\True Tổng các nghiệm của phương trình $f(x)=-\dfrac12$ trong khoảng $(0;\pi)$ bằng $\dfrac{3\pi}{2}$}
\loigiai{a) Sai: $f(-x)=-f(x)$, hàm số lẻ.

b) Sai: $D=\mathbb R$.

c) Sai: $\sin2x=2\sin x\cos x$.

d) Đúng: $\sin2x=-\dfrac12$ với $2x\in(0;2\pi)$ cho $2x=\dfrac{7\pi}6$ hoặc $2x=\dfrac{11\pi}6$, tức $x=\dfrac{7\pi}{12}$, $x=\dfrac{11\pi}{12}$; tổng bằng $\dfrac{3\pi}2$.}
\end{ex}
